\documentclass[12pt,a4paper,twoside,spanish]{article}      % Libro a 11 pt
\usepackage[utf8]{inputenc}
\usepackage[height=17.5cm,width=13.5cm]{geometry}
\usepackage[spanish]{babel}         % diccionario
\usepackage{epsfig}         % Graficos Postscript
\usepackage{tabularx}
\usepackage{sectsty}
\usepackage{float}
\usepackage{xcolor}

% Margenes 
\marginparwidth 0pt     \marginparsep 0pt
\topmargin   0pt        \textwidth   6.5in
\textheight 23cm

% Margen izq del txt en impares.
\setlength{\oddsidemargin}{.0001\textwidth}

% Margen izq del txt en pares.
\setlength{\evensidemargin}{-.04\textwidth}

% Anchura del texto
\setlength{\textwidth}{.99\textwidth}



%P rofundidad de enumeracion y tabla de contenidos
\setcounter{secnumdepth}{3}
\setcounter{tocdepth}{3}



% Nuevos Comandos

%Comandos para simplificar la escritura


\def\mc{\multicolumn}

% Comandos para poder utilizar raggedright en tablas

\newcommand{\PreserveBackslash}[1]{\let\temp=\\#1\let\\=\temp}
\let\PBS=\PreserveBackslash


%Cuerpo del documento


\begin{document}

    \def\chaptername{Capítulo}
    \def\tablename{Tabla}
    \def\listtablename{Índice de Tablas}
    \chapterfont{\LARGE\raggedleft}
    
    % DISEÑO DE LA PAGINA DEL TITULO

    \pagestyle{empty}
    
    \begin{titlepage}
        \setlength{\parindent}{0cm} \setlength{\parskip}{0cm}
        \newcommand{\HRule}{\rule{\linewidth}{1mm}}
        
        \vspace*{2cm}
        \HRule \\[0.5cm]
        \begin{center}
        % Letra lineal y negrita
        \textsf{\textbf{\large UN SISTEMA BASADO EN CONOCIMIENTO PARA \ldots(LA PLANIFIACIÖN Y ENTRENAMIENTO DEL KARATE)\\[1.5cm]
        Análisis de viabilidad e impacto. \\[0.25cm] Modelado del Contexto en CommonKADS \\[0.5cm]}}
        \HRule \vspace*{4cm}
        
        \textsf{\textbf{\normalsize Daniel Gago Martínez\\Miembro 2\\Miembro3\\Miembro 4 \\[1cm]
        Grupo de prácticas: 3.1\\[3.5cm]}}
        
        \textsf{\textbf{\small Desarrollo de Sistemas Inteligentes\\
        Universidade da Coruña \\ Curso 2025/26}}
        \end{center}
    \end{titlepage}
    
    \cleardoublepage
    \pagenumbering{Roman}
    \tableofcontents
    \cleardoublepage
    
    %numeros arábigos
    \pagenumbering{arabic} 
    %cabeceras
    \pagestyle{myheadings} \markboth{Título reducido} {Modelado Contextual en CommonKADS.}
    
    %indentaciones y espaciado entre párrafos
    \setlength{\parindent}{1,5cm} \setlength{\parskip}{0,7cm}
    
    \cleardoublepage
    
    \section{Análisis de Viabilidad: Modelado de la Organización.}
        \subsection{Formulario OM-1: contexto organizacional, problemas y soluciones.}

        
            \begin{table}[H]
                \scriptsize
                \begin{tabularx}{\textwidth}{|l|X|} \hline
                    \textbf{Modelo de Organización} & \textbf{Formulario OM-1: Problemas y Posibilidades de Mejora} \\ \hline\hline
                    
                    \textsc{Problemas y Oportunidades} & 
                        \begin{itemize}
                            \item El planificar el entrenamiento especifico para varios alumnos a la vez
                        \end{itemize}\\ \hline
                    
                    \textsc{Contexto Organizacional} & Nuestro sistema inteligente sera un producto propio que sera vendido a particulares:
                     
                    \begin{enumerate}
                        \item Planificar entrenamientos de karate para que los alumnos puedan entrenar de manera eficiente.
                        \item Que un estudiante se lesione y haya que rehacer todo su entrenamiento.
                        \item Un sensei que supervisa a los alumnos mientras realizan un entrenamiento en base a una planificación.
                        \item Valores que representan la organización: Los diferentes cinturones 
                    \end{enumerate}    \\ \hline
                    
                    \textsc{Soluciones} & Un sistema inteligente que permita realizar la planificación de varios entrenamientos particulares a a vez  \\ \hline
                \end{tabularx}
                  \label{tab.OM1}
            \end{table}
            
            
            \pagebreak
            
        \subsection{Formulario OM-2: descripción del área de interés de la organización.}

            

            \begin{table}[H]
            \scriptsize
                \begin{tabularx}{\textwidth}{|l|X|} \hline
                    \textbf{Modelo de Organización} & \textbf{Formulario OM-2: Aspectos Variables} \\ \hline\hline
                    
                    \textsc{Estructura} & 
                    Dojo de karate compuesto por:
                    \begin{itemize}
                        \item Dirección técnica, (Sensei)
                        \item Alumnos
                    \end{itemize} \\ \hline
                    
                    \textsc{Procesos} & 
                    Proceso principal:
                    \begin{itemize}
                        \item Planificación y gestión del entrenamiento
                    \end{itemize}
                    
                    Se descompone en tareas:
                    \begin{itemize}
                        \item Evaluación inicial del alumno
                        \item Definición de objetivos
                        \item Planificación del entrenamiento
                        \item Seguimiento del progreso
                        \item Evaluación de resultados
                    \end{itemize} \\ \hline
                    
                    \textsc{Personal} &  
                    Miembros implicados:
                    \begin{itemize}
                        \item \textbf{Sensei}: responsable técnico y docente
                        \item \textbf{Alumno}: destinatario del entrenamiento
                    \end{itemize} \\ \hline
                    
                    \textsc{Recursos} &  
                    Recursos utilizados:
                    \begin{enumerate}
                        \item Equipamiento y material: tatami, protecciones, cronómetros
                        \item Documentación y normativa: reglas y criterios federativos
                    \end{enumerate} \\ \hline
                    
                    \textsc{Conocimiento} &  
                    \begin{itemize}
                        \item \textbf{Criterios técnicos por cinturón}: dominio de técnicas y katas
                        \item \textbf {conocimiento relacion}: entre indicadores fisicos y ejercicios
                        \item \textbf{Reglas de diagnóstico}: evaluar el estado inicial del alumno
                        \item \textbf{Principios de planificación}: organización de sesiones y periodización
                        \item \textbf{Reglamento federativo}: normas y criterios de progresión
                        \item \textbf{Criterios de evaluación}: medir progreso y ajustar planes
                    \end{itemize}
                    (Véase OM-4 para detalle formal). \\ \hline
                    
                    \textsc{Cultura y Potencial} &  
                    El dojo sigue una cultura de disciplina, mejora continua y respeto a la jerarquía.  
                    Se valoran la experiencia del Sensei y la transmisión directa de conocimientos. 
                    
                    Experiencia interpersonal: comunicación y feedback entre Sensei y alumno. \\ \hline
                    
                \end{tabularx}
                \label{tab.OM2}
            \end{table}
            
            
            
            \pagebreak
        \subsection{Formulario OM-3: descomposición del proceso de negocio.}
    
        
        \begin{table}[H]
            \scriptsize
            \begin{tabularx}{\textwidth}{|p{0.2cm}|>{\raggedright}X|>{\raggedright}X|>{\raggedright}X|>{\raggedright}X|>{\raggedright}X|>{\PBS\raggedright}X|} \hline
            
                \multicolumn{3}{|l}{\textbf{Modelo de Organización}} & 
                \multicolumn{4}{|l|}{\textbf{Formulario OM-3: Descomposición de los Procesos}}\\ \hline\hline 
                
                \textsc{N\textordmasculine} &
                \textsc{Tarea} & 
                \textsc{Realiza\-da por} &  
                \textsc{¿Dónde?} & 
                \textsc{Recursos de Conocimiento} & 
                \textsc{¿In\-ten\-si\-va en Conocimiento?} &
                \textsc{Im\-por\-tan\-cia} \\ \hline
                
                1 & 
                Evaluación inicial del alumno & 
                Sensei & 
                Dojo& 
                conocimiento relacion entre indicadores fisicos y ejercicios & 
                Sí & 
                4 \\ \hline
                
            \end{tabularx}
            \label{tab.OM3-1}
        \end{table}

        \begin{table}[H]

        
        \scriptsize
            \begin{tabularx}{\textwidth}{|p{0.2cm}|>{\raggedright}X|>{\raggedright}X|>{\raggedright}X|>{\raggedright}X|>{\raggedright}X|>{\PBS\raggedright}X|} \hline
                \multicolumn{3}{|l}{\textbf{Modelo de Organización}} & 
                \multicolumn{4}{|l|}{\textbf{Formulario OM-3: Descomposición de los Procesos}}\\ \hline\hline 
                
                \textsc{N\textordmasculine} &
                \textsc{Tarea} & 
                \textsc{Realiza\-da por} &  
                \textsc{¿Dónde?} & 
                \textsc{Recursos de Conocimiento} & 
                \textsc{¿In\-ten\-si\-va en Conocimiento?} &
                \textsc{Im\-por\-tan\-cia} \\ \hline
                
                2 & 
                Definición de objetivos de entrenamiento & 
                Sensei& 
                Dojo& 
                Principios de planificación deportiva, progresión técnica, reglamento federativo (véase OM-4) & 
                Sí & 
                4 \\ \hline
                
            \end{tabularx}
            \label{tab.OM3-2}
        \end{table}

        \begin{table}[H]
        \scriptsize
            \begin{tabularx}{\textwidth}{|p{0.2cm}|>{\raggedright}X|>{\raggedright}X|>{\raggedright}X|>{\raggedright}X|>{\raggedright}X|>{\PBS\raggedright}X|} \hline      
                \multicolumn{3}{|l}{\textbf{Modelo de Organización}} & 
                \multicolumn{4}{|l|}{\textbf{Formulario OM-3: Descomposición de los Procesos}}\\ \hline\hline 
                
                \textsc{N\textordmasculine} &
                \textsc{Tarea} & 
                \textsc{Realiza\-da por} &  
                \textsc{¿Dónde?} & 
                \textsc{Recursos de Conocimiento} & 
                \textsc{¿In\-ten\-si\-va en Conocimiento?} &
                \textsc{Im\-por\-tan\-cia} \\ \hline
                
                3 & 
                Planificación del entrenamiento & 
                Sensei & 
                Dojo & 
                Principios de planificación, reglas de carga-descanso, modelos de periodización,Criterios técnicos por cinturón,Reglamento federativo (véase OM-4) & 
                Sí & 
                5 \\ \hline
            \end{tabularx}
        \label{tab.OM3-3}
        \end{table}

        \begin{table}[H]
        \scriptsize
            \begin{tabularx}{\textwidth}{|p{0.2cm}|>{\raggedright}X|>{\raggedright}X|>{\raggedright}X|>{\raggedright}X|>{\raggedright}X|>{\PBS\raggedright}X|} \hline
                
                \multicolumn{3}{|l}{\textbf{Modelo de Organización}} & 
                \multicolumn{4}{|l|}{\textbf{Formulario OM-3: Descomposición de los Procesos}}\\ \hline\hline 
                
                \textsc{N\textordmasculine} &
                \textsc{Tarea} & 
                \textsc{Realiza\-da por} &  
                \textsc{¿Dónde?} & 
                \textsc{Recursos de Conocimiento} & 
                \textsc{¿In\-ten\-si\-va en Conocimiento?} &
                \textsc{Im\-por\-tan\-cia} \\ \hline
                
                4 & 
                Seguimiento y monitorización & 
                Sensei& 
                Dojo& 
                Indicadores de rendimiento, umbrales de fatiga, reglas de control (véase OM-4) & 
                Sí & 
                4 \\ \hline
                
            \end{tabularx}
        \label{tab.OM3-4}
        \end{table}


        \begin{table}[H]
        \scriptsize
            \begin{tabularx}{\textwidth}{|p{0.2cm}|>{\raggedright}X|>{\raggedright}X|>{\raggedright}X|>{\raggedright}X|>{\raggedright}X|>{\PBS\raggedright}X|} \hline
            
                \multicolumn{3}{|l}{\textbf{Modelo de Organización}} & 
                \multicolumn{4}{|l|}{\textbf{Formulario OM-3: Descomposición de los Procesos}}\\ \hline\hline 
                
                \textsc{N\textordmasculine} &
                \textsc{Tarea} & 
                \textsc{Realiza\-da por} &  
                \textsc{¿Dónde?} & 
                \textsc{Recursos de Conocimiento} & 
                \textsc{¿In\-ten\-si\-va en Conocimiento?} &
                \textsc{Im\-por\-tan\-cia} \\ \hline
                
                5 & 
                Evaluación de resultados y reajuste & 
                Sensei& 
                Dojo& 
                Criterios de evaluación comparativa, métricas de progreso, reglas de adaptación (véase OM-4) & 
                Sí & 
                3 \\ \hline
                
            \end{tabularx}
        \label{tab.OM3-5}
        \end{table}
        
        
        \pagebreak

        
        \subsection{Formulario OM-4: activos de conocimiento.}
        
            \begin{table}[H]
                \scriptsize
                \begin{tabularx}{\textwidth}{|p{1.3cm}|p{1.3cm}|p{1.3cm}|X|X|X|X|} \hline
                
                    \multicolumn{3}{|l}{\textbf{Modelo de Organización}} & 
                    \multicolumn{4}{|l|}{\textbf{Formulario OM-4: Activos de Conocimiento}} \\ \hline\hline
                    
                    \textsc{Recurso de Conocimiento} & \textsc{Pertenece a} &  \textsc{Usado en} &  \textsc{¿Forma Correcta?} & \textsc{¿Lugar Correcto?} & \textsc{¿Tiempo Correcto?} & \textsc{¿Calidad Correcta?} \\ \hline
                    
                    Criterios técnicos por cinturón & Sensei & Evaluación inicial, Definición de objetivos, Planificación & Sí & Sí & Sí & Sí \\ \hline
                    
                \end{tabularx}
                \label{tab.OM4-criterios}
            \end{table}


            \begin{table}[H]
                \scriptsize
                \begin{tabularx}{\textwidth}{|p{1.3cm}|p{1.3cm}|p{1.3cm}|X|X|X|X|} \hline
                    \multicolumn{3}{|l}{\textbf{Modelo de Organización}} & 
                    \multicolumn{4}{|l|}{\textbf{Formulario OM-4: Activos de Conocimiento}} \\ \hline\hline
                    
                    \textsc{Recurso de Conocimiento} & \textsc{Pertenece a} &  \textsc{Usado en} &  \textsc{¿Forma Correcta?} & \textsc{¿Lugar Correcto?} & \textsc{¿Tiempo Correcto?} & \textsc{¿Calidad Correcta?} \\ \hline
                    
                    Indicadores físicos & Sensei & Evaluación inicial, Seguimiento del progreso & Sí & Sí & Sí & Sí \\ \hline
                    
                \end{tabularx}
                \label{tab.OM4-indicadores}
            \end{table}


            \begin{table}[H]
                \scriptsize
                \begin{tabularx}{\textwidth}{|p{1.3cm}|p{1.3cm}|p{1.3cm}|X|X|X|X|} \hline         
                    \multicolumn{3}{|l}{\textbf{Modelo de Organización}} & 
                    \multicolumn{4}{|l|}{\textbf{Formulario OM-4: Activos de Conocimiento}} \\ \hline\hline
                    
                    \textsc{Recurso de Conocimiento} & \textsc{Pertenece a} &  \textsc{Usado en} &  \textsc{¿Forma Correcta?} & \textsc{¿Lugar Correcto?} & \textsc{¿Tiempo Correcto?} & \textsc{¿Calidad Correcta?} \\ \hline
                    
                    Reglas de diagnóstico & Sensei & Evaluación inicial & Sí & Sí & Sí & Sí \\ \hline
                    
                \end{tabularx}
                \label{tab.OM4-reglas}
            \end{table}

            \begin{table}[H]
                \scriptsize
                \begin{tabularx}{\textwidth}{|p{1.3cm}|p{1.3cm}|p{1.3cm}|X|X|X|X|} \hline   
                    \multicolumn{3}{|l}{\textbf{Modelo de Organización}} & 
                    \multicolumn{4}{|l|}{\textbf{Formulario OM-4: Activos de Conocimiento}} \\ \hline\hline
                    
                    \textsc{Recurso de Conocimiento} & \textsc{Pertenece a} &  \textsc{Usado en} &  \textsc{¿Forma Correcta?} & \textsc{¿Lugar Correcto?} & \textsc{¿Tiempo Correcto?} & \textsc{¿Calidad Correcta?} \\ \hline
                    
                    Principios de planificación deportiva & Sensei & Definición de objetivos, Planificación & Sí & Sí & Sí & Sí \\ \hline
                    
                \end{tabularx}
                \label{tab.OM4-principios}
            \end{table}


            \begin{table}[H]
                \scriptsize
                \begin{tabularx}{\textwidth}{|p{1.3cm}|p{1.3cm}|p{1.3cm}|X|X|X|X|} \hline                  
                    \multicolumn{3}{|l}{\textbf{Modelo de Organización}} & 
                    \multicolumn{4}{|l|}{\textbf{Formulario OM-4: Activos de Conocimiento}} \\ \hline\hline
                    
                    \textsc{Recurso de Conocimiento} & \textsc{Pertenece a} &  \textsc{Usado en} &  \textsc{¿Forma Correcta?} & \textsc{¿Lugar Correcto?} & \textsc{¿Tiempo Correcto?} & \textsc{¿Calidad Correcta?} \\ \hline
                    
                    Reglamento federativo & Sensei & Definición de objetivos, Evaluación de resultados & Sí & Sí & Sí & Sí \\ \hline
                    
                \end{tabularx}
                \label{tab.OM4-reglamento}
            \end{table}


            \begin{table}[H]
                \scriptsize
                \begin{tabularx}{\textwidth}{|p{1.3cm}|p{1.3cm}|p{1.3cm}|X|X|X|X|} \hline
                    \multicolumn{3}{|l}{\textbf{Modelo de Organización}} & 
                    \multicolumn{4}{|l|}{\textbf{Formulario OM-4: Activos de Conocimiento}} \\ \hline\hline
                    
                    \textsc{Recurso de Conocimiento} & \textsc{Pertenece a} &  \textsc{Usado en} &  \textsc{¿Forma Correcta?} & \textsc{¿Lugar Correcto?} & \textsc{¿Tiempo Correcto?} & \textsc{¿Calidad Correcta?} \\ \hline
                    
                    Principios de planificación & Sensei & Planificación & Sí & Sí & Sí & Sí \\ \hline
                    
                \end{tabularx}
                \label{tab.OM4-heuristicas}
            \end{table}


            \begin{table}[H]
                \scriptsize
                \begin{tabularx}{\textwidth}{|p{1.3cm}|p{1.3cm}|p{1.3cm}|X|X|X|X|} \hline   
                    \multicolumn{3}{|l}{\textbf{Modelo de Organización}} & 
                    \multicolumn{4}{|l|}{\textbf{Formulario OM-4: Activos de Conocimiento}} \\ \hline\hline
                    
                    \textsc{Recurso de Conocimiento} & \textsc{Pertenece a} &  \textsc{Usado en} &  \textsc{¿Forma Correcta?} & \textsc{¿Lugar Correcto?} & \textsc{¿Tiempo Correcto?} & \textsc{¿Calidad Correcta?} \\ \hline
                    
                    Umbrales de fatiga y rendimiento & Sensei & Seguimiento, Evaluación de resultados & Sí & Sí & Sí & Sí \\ \hline           
                \end{tabularx}
                \label{tab.OM4-umbrales}
            \end{table}



            \begin{table}[H]
                \scriptsize
                \begin{tabularx}{\textwidth}{|p{1.3cm}|p{1.3cm}|p{1.3cm}|X|X|X|X|} \hline           
                    \multicolumn{3}{|l}{\textbf{Modelo de Organización}} & 
                    \multicolumn{4}{|l|}{\textbf{Formulario OM-4: Activos de Conocimiento}} \\ \hline\hline
                    
                    \textsc{Recurso de Conocimiento} & \textsc{Pertenece a} &  \textsc{Usado en} &  \textsc{¿Forma Correcta?} & \textsc{¿Lugar Correcto?} & \textsc{¿Tiempo Correcto?} & \textsc{¿Calidad Correcta?} \\ \hline
                    
                    Criterios de evaluación comparativa & Sensei & Evaluación de resultados & Sí & Sí & Sí & Sí \\ \hline                
                \end{tabularx}
                \label{tab.OM4-criterios-eval}
            \end{table}

            \pagebreak
            
        \subsection{Formulario OM-5: Análisis de viabilidad.}
        
        \begin{table}[H]
            \scriptsize
            \begin{tabularx}{\textwidth}{|l|X|} \hline

                \textbf{Modelo de Organización} & \textbf{Formulario OM-5: Elementos del Documento de Viabilidad}\\ \hline\hline
                
                \textsc{Viabilidad Empresarial} & Se deben responder las siguientes preguntas para cada problema y posibilidad e incluir la solución propuesta:
                
                \begin{enumerate}
                    \item¿Cuáles son los beneficios esperados para la organización con la adopción de la solución? Se deben identificar los beneficios tanto económicos como intangibles.
                    
                    \item¿Cuánto es el valor añadido esperado?
                    
                    \item¿Cuál es el coste esperado de la solución? 
                    
                    \item¿Y comparado con otras soluciones alternativas?
                    
                    \item¿Se requieren cambios organizacionales?
                    
                    \item¿Hasta qué punto están implicados los riesgos e incertidumbres de la empresa en la solución propuesta?
                \end{enumerate}
                \\ \hline
            \end{tabularx}
            \caption{Formulario OM-5 (parte I).}
              \label{tab.OM5_1}
        \end{table}
        
        \begin{table}[H]
            \scriptsize
            \begin{tabularx}{\textwidth}{|l|X|} \hline
            
                \textbf{Modelo de Organización} & \textbf{Formulario OM-5: Elementos del Documento de Viabilidad}\\ \hline\hline
                
                \textsc{Viabilidad Técnica} & Se deben responder las siguientes preguntas para cada problema y posibilidad junto
                a la solución propuesta:
                
                \begin{enumerate}
                    \item ¿Qué complejidad tiene la tarea que desarrollará el SBC desde el punto de vista del conocimiento a almacenar y los procesos de razonamiento a realizar? ¿Son las técnicas y métodos actuales adecuados?
                    
                    \item ¿Existen aspectos críticos relativos al tiempo, calidad, recursos necesarios, etc.? ¿Qué se va a hacer al respecto?
                    
                    \item ¿Están claros los criterios bajo los cuales vamos a determinar el éxito del proyecto? ¿Cómo vamos a comprobar la calidad y la validez del sistema?
                    
                    \item ¿Cómo son de complejas las interfaces con el usuario que se necesitan? ¿Son las técnicas y métodos actuales adecuados?
                    
                    \item ¿Cómo es de compleja la interacción con otros sistemas de información y otros recursos potenciales (integración de sistemas, interoperatividad)? ¿Son las técnicas y métodos actuales adecuados?
                    
                    \item ¿Existen riesgos e incertidumbres tecnológicos adicionales?
                \end{enumerate}   \\ \hline
                
                
                \textsc{Viabilidad del Proyecto} & Se deben responder las siguientes preguntas para cada problema y posibilidad junto
                a la solución propuesta:
                
                \begin{enumerate}
                    \item ¿Existe el compromiso adecuado por parte del personal (gestores, expertos, usuarios, clientes, equipo del proyecto) para llevar a cabo los siguientes pasos del proyecto?
                    
                    \item ¿Están disponibles los recursos necesarios (en tiempo, presupuesto, equipamiento, personal)?
                    
                    \item ¿Están disponibles el conocimiento y las capacidades requeridas?
                    
                    \item ¿Son realistas las expectativas sobre el proyecto y sus resultados?
                    
                    \item ¿Es adecuada la organización del proyecto y las comunicaciones externas o internas?
                    
                    \item ¿Existen riesgos e incertidumbres adicionales relativos al proyecto?
                \end{enumerate}
                
                \\ \hline
                \textsc{Acciones Propuestas} & Esta parte del documento de viabilidad está directamente relacionada con los acuerdos de gestión y la toma de decisiones. Valora e integra en un plan de actuación concreto los resultados obtenidos en el análisis previo. Se va completando a lo largo de todo el ciclo de vida del proyecto (en espiral). Para esta primera etapa, tan solo es necesario responder a las preguntas 1 y 2.
                
                \begin{enumerate}
                    \item ¿Qué área de actuación se recomienda entre las áreas de problemas/oportunidades encontradas y descritas en OM-1?
                    
                    \item ¿Cuál es la solución recomendada para el área elegida y qué tareas de OM-3 deberían automatizarse?
                    
                    \item ¿Cuáles son los resultados, costes y beneficios esperados?
                    
                    \item ¿Qué acciones deben de tomar dentro del proyecto para alcanzar dicha solución?
                    
                    \item Si las circunstancias que rodean a la organización cambian (tanto externas como internas), ¿en qué condiciones es apropiado reconsiderar las decisiones tomadas?
                \end{enumerate}
                \\ \hline
            \end{tabularx}
            \caption{Formulario OM-5 (parte II).}
              \label{tab.OM5_2}
        \end{table}
        
        
        \pagebreak
        
    \section{Análisis de Impactos y Mejoras: Modelado de las Tareas y los Agentes.}
    
        \subsection{Formulario TM-1: análisis de tareas.}



        \begin{table}[H]
            \scriptsize
            \begin{tabularx}{\textwidth}{|l|X|} \hline       
                \textbf{Modelo de Tareas} & \textbf{Formulario TM-1: Análisis de Tareas} \\ \hline\hline
                
                \textsc{Tarea} & T3 – Planificación del entrenamiento (Ref. OM-3) \\ \hline
                
                \textsc{Organización}  & Proceso central del dojo bajo responsabilidad técnica del Sensei. \\ \hline
                
                \textsc{Objetivo y valor} & Diseñar sesiones estructuradas que permitan alcanzar los objetivos definidos. Aporta organización y eficiencia. \\ \hline
                
                \textsc{Dependencia y Flujos} 
                & \textit{1. Tareas precedentes:} T2 – Definición de objetivos.\\
                & \textit{2. Tareas que le siguen:} T4 – Seguimiento del progreso. \\ \hline
                
                \textsc{Objetos manipulados} 
                & 1. Entrada: objetivos definidos.\\
                & 2. Salida: plan de entrenamiento estructurado.\\
                & 3. Internos: heurísticas y principios de planificación (OM-4). \\ \hline
                
                \textsc{Tiempo y control} 
                & 1. Frecuencia: periódica (mensual o por ciclo).\\
                & 2. Control: determina ejecución de sesiones.\\
                & 3. Restricciones: equilibrio entre carga y recuperación. \\ \hline
                
                \textsc{Agentes} & Sensei. \\ \hline
                
                \textsc{Conocimiento y Capacidad} & Uso intensivo de heurísticas y principios deportivos (OM-4). Tarea altamente intensiva en conocimiento. \\ \hline
                
                \textsc{Recursos} & Tiempo del Sensei, material del dojo. \\ \hline
                
                \textsc{Calidad y eficiencia} & Coherencia del plan y progresión técnica observada. \\ \hline
            \end{tabularx}
            \label{tab.TM1-T3}
        \end{table}

        
        
        \pagebreak
        
    \subsection{Formulario TM-2: análisis de los cuellos de botella del conocimiento.}
    
        \begin{table}[H]
            \scriptsize
            \begin{tabularx}{\textwidth}{|p{5cm}|>{\PBS\raggedright}p{0.8cm}|X|} \hline
                \textbf{Modelo de Tareas} & \multicolumn{2}{l|}{\textbf{Formulario TM-2: Elemento de Co\-no\-ci\-mien\-to}} \\ \hline\hline
                
                \textsc{Nombre} &  \multicolumn{2}{l|}{Elemento de conocimiento \emph{(Referencia a OM-3)}}\\ \hline
                
                \textsc{Poseído por} &  \multicolumn{2}{X|}{Agente \emph{(Referencia a OM-4)}}\\ \hline
                
                \textsc{Usado en} &  \multicolumn{2}{l|}{Nombre de la tarea e identificador \emph{(Referencia a OM-3)}}\\ \hline
                
                \textsc{Dominio} &  \multicolumn{2}{p{7.5cm}|}{Dominio más amplio en el que se encuentra el conocimiento (campo de la especialidad,
                disciplina, rama de la ciencia o ingeniería, etc.).}\\ \hline
                
                \textbf{Naturaleza del conocimiento} & \emph{(Sí/No)} &
                \textbf{¿Supone un cuello de botella?¿Debe ser mejorado?}\\ \hline 
                
                Formal, riguroso & & \\ \hline 
                
                Empírico, cuantitativo& & \\ \hline
                
                Heurístico, sentido común& & \\ \hline
                
                Altamente especializado, específico del dominio& & \\ \hline 
                
                Basado en la experiencia& & \\ \hline 
                
                Basado en la acción$^1$& & \\ \hline
                
                Incompleto& & \\\hline 
                
                Incierto, puede ser incorrecto & & \\ \hline 
                
                Cambia con rapidez& & \\ \hline 
                
                Difícil de verificar& & \\ \hline 
                
                Tácito, difícil de transferir& & \\ \hline 
                
                \textbf{Forma del conocimiento} & &\\ \hline
                
                Mental& & \\ \hline
                
                Papel& & \\ \hline 
                
                Electrónica& & \\ \hline
                
                Habilidades& & \\ \hline 
                
                Otros& & \\ \hline 
                
                \textbf{Disponibilidad del Conocimiento} &  &\\ \hline 
                
                Limitaciones en tiempo& & \\ \hline
                
                Limitaciones en espacio& & \\ \hline 
                
                Limitaciones de acceso& & \\ \hline Limitaciones de calidad& & 
                
                \\ \hline Limitaciones de forma& & \\ \hline
                
            \end{tabularx}
              \label{tab.TM2}
        \end{table}
        
        {\scriptsize \noindent $^1$ Se refiere a conocimiento que se
        adquiere con la repetición de actividades físicas tales como
        conducir, encestar, etc.}
        
        \pagebreak
        
    \subsection{Formulario AM-1: descripción de los agentes.}
        

        % =================== AGENTE 1: SENSEI ===================

        \begin{table}[H]
            \scriptsize
            \begin{tabularx}{\textwidth}{|l|X|} \hline
                \textbf{Modelo de Agentes} & \textbf{Formulario AM-1: Agentes} \\ \hline\hline
                
                \textsc{Nombre} &  Sensei \\ \hline
                
                \textsc{Organización} &  Dirección técnica, humano \emph{(referencia a OM-2)}\\ \hline
                
                \textsc{Implicado en} & T2, definición de objetivos; T3, planificación de entrenamiento; T4, seguimiento del progreso \emph{(referencia a TM-1)} \\ \hline
                
                \textsc{Se comunica con} &  Alumnos, Sistema inteligente  \\ \hline
                
                \textsc{Conocimiento} &  Elementos de conocimiento que el agente posee \emph{(referencia a TM-2)} \\ \hline
                
                \textsc{Otras Competencias} &  Experiencia técnica, liderazgo \\ \hline
                
                \textsc{Responsabilidades y restricciones} &  Supervisar la correcta progresión de los alumnos, respetar las normas del dojo, garantizar la seguridad de los alumnos\\ \hline
            \end{tabularx}
             \label{tab.AM1}
        \end{table}
        

        \pagebreak


        % =================== AGENTE 2: ALUMNO ===================

        \begin{table}[H]
            \scriptsize
            \begin{tabularx}{\textwidth}{|l|X|} \hline
                \textbf{Modelo de Agentes} & \textbf{Formulario AM-1: Agentes} \\ \hline\hline
                
                \textsc{Nombre} &  Alumno \\ \hline
                
                \textsc{Organización} &  Miembro del dojo, humano \emph{(referencia a OM-2)}\\ \hline
                
                \textsc{Implicado en} & Ejecutar el plan de entrenamiento, registrar el progreso (relacionado con T3 y T4) \emph{(referencia a TM-1)} \\ \hline
                
                \textsc{Se comunica con} &  Sensei, alumnos  \\ \hline
                
                \textsc{Conocimiento} &  Elementos de conocimiento que el agente posee \emph{(referencia a TM-2)} \\ \hline
                
                \textsc{Otras Competencias} &  Autogestión, capacidad de aprendizaje, disciplina \\ \hline
                
                \textsc{Responsabilidades y restricciones} &  Seguir las indicaciones del sensei, cumplir el plan, informar de lesiones\\ \hline
            \end{tabularx}
             \label{tab.AM1}
        \end{table}
        
        
        \pagebreak


        % =================== AGENTE 3: SISTEMA DE PLANIFICACIÓN DE ENTRENAMIENTO ===================

        \begin{table}[H]
            \scriptsize
            \begin{tabularx}{\textwidth}{|l|X|} \hline
                \textbf{Modelo de Agentes} & \textbf{Formulario AM-1: Agentes} \\ \hline\hline
                
                \textsc{Nombre} &  Sistema de planificaciónde entrenamiento \\ \hline
                
                \textsc{Organización} &  Departamento de planificación, sistema inteligente \emph{(referencia a OM-2)}\\ \hline
                
                \textsc{Implicado en} & T3, planificación de entrenamiento \emph{(referencia a TM-1)} \\ \hline
                
                \textsc{Se comunica con} &  Sensei, alumnos  \\ \hline
                
                \textsc{Conocimiento} &  Elementos de conocimiento que el agente posee \emph{(referencia a TM-2)} \\ \hline
                
                \textsc{Otras Competencias} &  Recordatorios, reportes de progreso \\ \hline
                
                \textsc{Responsabilidades y restricciones} &  Garantizar planes adecuados al nivel del alumno, respetar limitaciones del alumno\\ \hline
            \end{tabularx}
             \label{tab.AM1}
        \end{table}
        
        
        \pagebreak
        
    \section{Conclusiones del análisis de viabilidad e impacto.}
    
        \subsection{Formulario OTA-1: recomendaciones y acciones.}
               
            \vspace*{1cm}
                    
            \begin{table}[H]
                \scriptsize
                \begin{tabularx}{\textwidth}{|p{4cm}|X|} \hline
                    \textbf{Modelos de Organización, Tareas y Agentes} & \textbf{Formulario OTA-1: Documento sobre Impactos y Mejoras}\\ \hline\hline 
                    
                    \textsc{Impactos y Cambios en la Organización} & Describir los impactos y los cambios que el SBC traerá a la organización, comparándola con la estructura organizativa actual \emph{(descrita en OM-2)}, y con cómo será dicha organización en el futuro. Esta descripción se realizará de forma global para los seis aspectos variables descritos en OM-2. \\ \hline
                    
                    \textsc{Impactos y Cambios en Tareas y Agentes} & Describir los impactos y los cambios que el SBC introducirá en los agentes y las tareas, comparándolos con la situación actual \emph{(descrita en TM-1/2 y AM-1)}. Es importante que no sólo se considere al personal implicado en una determinada tarea, sino también a otros agentes y poseedores de conocimiento (gestores, usuarios, clientes): \\
                    
                    & 1. Cambios en la estructura de las tareas (flujo, dependencias, objetos manipulados, temporización, control). \\
                    & 2. Cambios en los recursos necesarios. \\
                    & 3. Criterios para evaluar la calidad y el rendimiento.\\
                    & 4. Cambios en la plantilla y en los agentes implicados.\\
                    & 5. Cambios en las posiciones, responsabilidades, autoridad y limitaciones de los
                    individuos implicados en la ejecución de la tarea.\\
                    & 6. Cambios en el conocimiento y capacidad requeridos.\\
                    & 7. Cambios en los canales de comunicación.\\ \hline
                    
                    \textsc{Actitudes y Compromisos} &  Analizar cómo reaccionarán a los cambios introducidos los individuos y el personal cualificado
                    involucrado, y decidir si existe la base suficiente para llevar a cabo dichos cambios. \\ \hline
                    
                    \textsc{Acciones Propuestas} &  Apartado directamente relacionado con los acuerdos de gestión y la toma de decisiones. Los resultados obtenidos con el análisis previo se valoran e integran en un plan de actuación concreto: \\
                    
                    & \textit{1. Mejoras:} ¿Cuáles son los cambios recomendados (organización, tareas, personal, sistemas\ldots)?\\
                    & \textit{2. Medidas adicionales:} ¿Qué medidas hay que tomar para facilitar dichos cambios? (formación, facilidades\ldots)\\
                    & \textit{3. Acciones del proyecto:} ¿Cuál es la siguiente acción a realizar dentro del proyecto respecto a la solución basada en el conocimiento asumida?\\
                    & \textit{4. Resultados, costes y beneficios esperados} \emph{(referencia a OM-5)}.\\
                    & 5. Si las circunstancias que rodean a la organización cambian (tanto externas como internas), ¿en qué condiciones es apropiado reconsiderar las decisiones tomadas? \\ \hline
                \end{tabularx}
                  \label{tab.OTA1}
            \end{table}
            
        
\end{document}
